La completación cúbica y el transporte aperiódico precisan dos sentidos
de conservación que se articulan en el desarrollo. La primera conserva
la unidad bajo proyección de medidas normalizadas y distingue el cubo
inicial, la completación perforada de \(999\) elementos y la restitución
del ápice. El segundo conserva el registro de trayectoria durante el
retorno de fase. La monodromía y el modelo matemático de cristal temporal
aperiódico expresan este segundo mecanismo mediante aplicaciones y
observables, no sólo mediante una imagen helicoidal. La exponencial
trítica, por su parte, unifica las realizaciones elíptica, parabólica e
hiperbólica sin suprimir sus relaciones cuadráticas respectivas.

La constante de Planck reducida, \(\hbar\), ocupa una posición explícita
en esta arquitectura. El cociente determinantal de orden \(16\) actúa
en la norma que determina la década \(-34\); el carácter \(H_5\) y su
continuación regional determinan las dos secciones de acción cuyo
cociente entra en la coordenada electrónica. La comparación de la
sección anterior con \(h/(2\pi)\) muestra una diferencia relativa del
orden de \(10^{-15}\), y la sección retornada conserva la compensación
que utiliza el electrón. La escala de acción enlaza así la norma
determinantal y el carácter regional con el cociente electrónico: la
compensación pertenece a la construcción exacta, mientras que la proximidad
de la sección anterior a \(h/(2\pi)\) es una comparación numérica y no
una igualdad exacta.
